
\documentclass[11pt]{article}

\usepackage[margin=1in]{geometry}
\usepackage{graphicx}
\usepackage{booktabs}
\usepackage{amsmath}
\usepackage{float}
\usepackage{hyperref}

\title{Data Exploration, Wrangling, and Feature Engineering Using the ADNI Dataset}
\author{Melvin Gates Jr. \\
Artificial Intelligence / Data Science}
\date{January 23, 2026}

\begin{document}
\maketitle

\section{Project Overview}
In real-world machine learning and data science applications, model performance is strongly influenced by data quality, appropriate feature representation, and the dimensionality of the dataset. This project focuses on developing professional data preparation and feature engineering skills using clinical data from the Alzheimer’s Disease Neuroimaging Initiative (ADNI).

The project involves data exploration, data cleaning, feature engineering, and dimensionality reduction. All methodological decisions are justified based on statistical reasoning, clinical context, and potential downstream machine learning impact.

\section{Dataset Description}
The dataset used in this project was obtained from the ADNI repository, specifically the \texttt{TADPOLE\_D3.xlsx} file. The analysis utilized the \texttt{in} sheet, which contains demographic, clinical, and imaging-derived features.

After data cleaning, the dataset contained 896 observations and 365 features. Each observation represents a single subject at one clinical visit, corresponding to a cross-sectional snapshot rather than longitudinal measurements.

\section{Data Cleaning}

\subsection{Missing Data}
Missing values were identified for each feature by computing both the total count and the percentage of missing observations. Several features exhibited 100\% or near-100\% missing values and were removed because they contained no meaningful information and would introduce noise if retained.

\subsection{Incorrect or Inconsistent Data}
Diagnosis labels contained a mixture of static diagnostic categories (NL, MCI, Dementia) and transition-based labels (e.g., MCI to Dementia). Transition labels were mapped to their corresponding current diagnostic category to ensure consistency with the unit of analysis.

Additionally, numerical outliers were detected in the WholeBrain feature using summary statistics and boxplot visualization. An interquartile range–based capping approach was applied to reduce the influence of extreme values while preserving the sample size.

\subsection{Unnecessary Data}
Identifier variables such as RID, VISCODE, and EXAMDATE were removed because they do not provide predictive clinical information and may introduce information leakage. Features with extremely high missingness were also removed. No near-zero variance features were identified after cleaning.

\section{Feature Engineering}

\subsection{Feature Selection}
Three features were selected based on clinical relevance:
\begin{itemize}
\item AGE: A primary demographic risk factor for Alzheimer’s disease.
\item DX: Clinical diagnosis representing disease stage.
\item WholeBrain: A neuroimaging biomarker associated with neurodegeneration.
\end{itemize}

\subsection{Encoding}
Categorical features require numerical encoding for machine learning. Diagnosis was encoded using ordinal encoding to preserve its natural ordering (NL $\rightarrow$ MCI $\rightarrow$ Dementia) while minimizing dimensionality.

\subsection{Scaling}
Scaling was required due to differing feature magnitudes. Standardization (z-score scaling) was selected to ensure features have zero mean and unit variance, making it suitable for dimensionality reduction techniques such as PCA.

\section{Curse of Dimensionality and Dimensionality Reduction}

\subsection{Curse of Dimensionality}
The ADNI dataset is susceptible to the curse of dimensionality due to its large number of features relative to observations. High dimensionality increases data sparsity, computational cost, and the risk of overfitting.

\subsection{Dimensionality Reduction Using PCA}
Principal Component Analysis (PCA) was applied to the cleaned and standardized numerical features. PCA transforms correlated variables into orthogonal components that capture maximum variance. Reducing the dataset to two components enabled visualization and mitigated high-dimensional challenges.

\section{Conclusion}
This project demonstrated a complete data preparation and feature engineering pipeline for a real-world clinical dataset. Through systematic cleaning, encoding, scaling, and dimensionality reduction, the dataset was prepared for downstream machine learning analysis in a reproducible and interpretable manner.

\end{document}
